\section{Recovery of a prescribed resummed branch amplitude}
\label{sec:branch}
Fixed-order logarithmic polynomials do not cover general resummed celestial
amplitudes. In particular, branch structures from renormalization-group
resummation are already established in the literature \cite{gu}.
We give a separate recovery result for one specified logarithmic template.
Its functional form and parameters are assumptions of the theorem.

Let $\mathcal H$ be a complex Hilbert space (the scalar case is
$\mathcal H=\mathbb C$). Fix $a,\alpha,M,\lambda,\sigma>0$ and
$0<\gamma<1$. In a fixed numerical energy convention define
\begin{equation}\label{eq:branchtemplate}
 K(t)=(1+\lambda t)^{-\gamma}\quad(t\ge0),\qquad
 \phi(x)=K(\log(1/x))\mathbf1_{\{0<x\le1\}}.
\end{equation}
Consider strongly measurable kernels
\begin{equation}\label{eq:branchclass}
 \begin{gathered}
 f(x)=b\phi(x)+R(x),\quad b\in\mathcal H,\qquad
 \|R(x)\|_{\mathcal H}\le Mx^\alpha\quad(x>0),\\
 Q_a(f)=\int_0^\infty x^{a-1}\|f(x)\|_{\mathcal H}\,dx<\infty.
 \end{gathered}
\end{equation}
The remainder bound is global, including the ultraviolet region. It is
considerably stronger than an unspecified asymptotic expansion. The amplitude
$b$ is unique: if two decompositions satisfy this bound, divide their
difference by $\phi(x)$ and let $x\downarrow0$.
Set $G(z)=\int_0^\infty x^{z-1}f(x)\,dx$ on its convergence region.
The known template has transform
\begin{equation}\label{eq:branchtransform}
 H_\gamma(z)=\int_0^\infty e^{-zt}(1+\lambda t)^{-\gamma}\,dt
 =\lambda^{-\gamma}e^{z/\lambda}z^{\gamma-1}
   \Gamma(1-\gamma,z/\lambda),\qquad \Re z>0.
\end{equation}
Here $\Gamma(s,u)$ is the upper incomplete gamma function and the power uses
the principal branch. Substitution $v=t+1/\lambda$ proves the identity.
The remainder transform is holomorphic on $-\alpha<\Re z<a$:
the infrared part is controlled by $Mx^\alpha$, and the ultraviolet part
by $Q_a(f)$. Consequently
\begin{equation}\label{eq:branchcoefficient}
 G(z)=bH_\gamma(z)+G_R(z),\qquad
 \lim_{z\downarrow0}z^{1-\gamma}G(z)
 =b\lambda^{-\gamma}\Gamma(1-\gamma).
\end{equation}
The limit is along the positive real axis. The same leading term holds in
closed sectors of the slit neighborhood of zero. The target $b$ therefore
differs from the leading fractional-branch coefficient by a known nonzero
factor, rather than being a residue of a meromorphic pole.

For bounded-noise data $Y(\nu)=G(a+i\nu)+e(\nu)$ with
$\sup_\nu\|e(\nu)\|_{\mathcal H}\le\eta$, use the Gaussian functional
\begin{equation}\label{eq:branchgaussian}
 E_{h,\sigma}(Y)=\frac1{2\pi}\int_{\mathbb R}
 Y(\nu)h^{-a-i\nu}\exp\bigl(\sigma^2(a+i\nu)^2/2\bigr)\,d\nu.
\end{equation}
All vector-valued integrals are Bochner integrals. Put
\[
 A_{\alpha,\sigma}=e^{\alpha^2\sigma^2/2},\qquad
 B_{a,\sigma}=\frac{e^{a^2\sigma^2/2}}{\sqrt{2\pi}\sigma},
 \qquad L=-\log h.
\]
Define the exact, noiseless response of the prescribed template by
\begin{equation}\label{eq:branchresponse}
 \Phi_h=\mathbb E\,\phi(he^{\sigma Z})
 =\frac1{\sqrt{2\pi}\sigma}\int_0^\infty
 e^{-(t-L)^2/(2\sigma^2)}K(t)\,dt>0,
 \qquad Z\sim N(0,1).
\end{equation}
The normalization is computed from known $\lambda,\gamma$, not estimated
from noisy data.

\begin{theorem}[Known-branch recovery with a matching rate]
\label{thm:branchrecovery}
On the class \eqref{eq:branchclass}, the estimator
$\widehat b_h=E_{h,\sigma}(Y)/\Phi_h$ satisfies
\begin{equation}\label{eq:brancherror}
 \|\widehat b_h-b\|_{\mathcal H}
 \le\frac{A_{\alpha,\sigma}Mh^\alpha+
                B_{a,\sigma}\eta h^{-a}}{\Phi_h}.
\end{equation}
Let $p_0=\mathbb P(|Z|\le1)=\operatorname{erf}(1/\sqrt2)$.
For $L\ge\sigma$,
\begin{equation}\label{eq:branchresponselower}
 \Phi_h\ge p_0[1+\lambda(L+\sigma)]^{-\gamma}.
\end{equation}
With $\theta=\alpha/(a+\alpha)$ and
$S_\eta=M^{a/(a+\alpha)}\eta^\theta$, choosing
$h=(\eta/M)^{1/(a+\alpha)}$ gives
\begin{equation}\label{eq:branchupperrate}
 \|\widehat b_h-b\|_{\mathcal H}
 \le\frac{A_{\alpha,\sigma}+B_{a,\sigma}}{p_0}
 S_\eta
 \left[1+\lambda\left(\frac{\log(M/\eta)}{a+\alpha}
                         +\sigma\right)\right]^\gamma
\end{equation}
whenever this choice has $L\ge\sigma$.

For the scalar class, and also for any nonzero Hilbert space, let
$\mathcal R_{\mathrm{br}}(\eta)$ denote the minimax error for $b$.
Set
\[
 L_\eta=\frac{\log(M/(2a\eta))}{a+\alpha}.
\]
If $L_\eta\ge0$ and
$e^{-\alpha L_\eta}(1+\lambda L_\eta)^\gamma\le1$, then
\begin{equation}\label{eq:branchlowerrate}
 \mathcal R_{\mathrm{br}}(\eta)
 \ge\tfrac12(2a)^\theta S_\eta
                  (1+\lambda L_\eta)^\gamma.
\end{equation}
Thus, for fixed positive parameters and $\eta/M\downarrow0$,
\begin{equation}\label{eq:branchminimax}
 \mathcal R_{\mathrm{br}}(\eta)
 =\Theta\!\left(
 M^{a/(a+\alpha)}\eta^{\alpha/(a+\alpha)}
 [\log(M/\eta)]^\gamma\right).
\end{equation}
The notation compares rates with parameter-dependent constants. It does not
assert equality of the upper and lower constants.
\end{theorem}
\begin{proof}
The Gaussian inversion identity, justified by $Q_a(f)<\infty$ and the
integrable frequency kernel, gives
$E_{h,\sigma}(G)=\mathbb E f(he^{\sigma Z})$.
The template component is exactly $b\Phi_h$. The global remainder bound
therefore gives bias at most
$M\mathbb E(he^{\sigma Z})^\alpha=A_{\alpha,\sigma}Mh^\alpha$.
The frequency kernel in \eqref{eq:branchgaussian} has $L^1$ norm
$B_{a,\sigma}h^{-a}$, proving \eqref{eq:brancherror} after division by
$\Phi_h$. On $|Z|\le1$ and $L\ge\sigma$, the logarithmic energy
$L-\sigma Z$ lies in $[0,L+\sigma]$. Since $K$ is decreasing, the template
is at least $K(L+\sigma)$ on this event. This proves
\eqref{eq:branchresponselower}. The stated choice of $h$ makes both
$Mh^\alpha$ and $\eta h^{-a}$ equal to $S_\eta$, yielding the upper bound.

For the lower bound, take a unit vector $v\in\mathcal H$ and abbreviate
$L=L_\eta$. Compare zero with the positive scalar profile times $v$,
\[
 f_L(x)=b_L\phi(x)\mathbf1_{\{x\le e^{-L}\}}v,
 \qquad b_L=M e^{-\alpha L}/K(L).
\]
Its template amplitude is $b_Lv$. The remainder vanishes outside
$e^{-L}<x\le1$. On this interval put $t=\log(1/x)\in[0,L]$ and
$F(t)=e^{-\alpha t}/K(t)$. The function
$\log F(t)=-\alpha t+\gamma\log(1+\lambda t)$ is concave. Hence
$F(t)\ge\min\{F(0),F(L)\}=F(L)$ under the stated condition.
It follows that
\[
 \|R_L(x)\|x^{-\alpha}
 =b_LK(t)e^{\alpha t}=b_L/F(t)\le M.
\]
The Mellin integral converges, and positivity gives its exact line norm
as its value at zero frequency. Monotonicity of $K$ then gives
\[
 \sup_\nu\|G_L(a+i\nu)\|
 =b_L\int_L^\infty e^{-at}K(t)\,dt
 \le b_LK(L)e^{-aL}/a
 =M e^{-(a+\alpha)L}/a=2\eta.
\]
The observation $G_L/2$ is within $\eta$ of both exact data sets.
Their amplitudes differ by $b_L$, so every estimator has worst-case
error at least $b_L/2$, which is \eqref{eq:branchlowerrate}.
The conditions on $L_\eta$ hold for all sufficiently small $\eta/M$;
comparison with \eqref{eq:branchupperrate} proves \eqref{eq:branchminimax}.
Multiplying either bound by $\lambda^{-\gamma}\Gamma(1-\gamma)$ gives
the corresponding bounds for the leading fractional-branch coefficient.
\end{proof}

\begin{corollary}[Normalized finite-band recovery]
\label{cor:branchband}
If $Q_a(f)\le Q$ is also known, restrict the integral in
\eqref{eq:branchgaussian} to $|\nu|\le\Omega$ but retain the exact
template normalization $\Phi_h$. Then
\[
 \|\widehat b_{h,\Omega}-b\|
 \le\frac{A_{\alpha,\sigma}Mh^\alpha+
 B_{a,\sigma}h^{-a}
 [\eta+Q\operatorname{erfc}(\sigma\Omega/\sqrt2)]}{\Phi_h}.
\]
This assumes continuum-band observations. Sampling and quadrature error
need their own bounds.
\end{corollary}
\begin{proof}
Use $\|G(a+i\nu)\|\le Q$ on the omitted frequencies and integrate the
absolute Gaussian tail. The remaining bias and noise bounds are those in
Theorem~\ref{thm:branchrecovery}.
\end{proof}

The exponent and template shape are known inputs here. If $\gamma$ is
unknown and the amplitude may vanish, $b=0$, $R=0$ gives identical exact
data for every exponent, so $\gamma$ is not identifiable on that class.
Even if $b$ is required to be positive but has no uniform lower bound,
uniformly stable exponent recovery is impossible: for any two fixed
exponents, $\|bH_{\gamma_1}-bH_{\gamma_2}\|_{a,\infty}\le2b/a$,
because $0<K(t)\le1$. Taking $0<b\le a\eta$ supplies a common
noise-ball observation while leaving the two exponents separated.
Joint parameter estimation with a nonzero amplitude lower bound is a
different problem.

The normalization removes the avoidable $O(1/L)$ approximation error that
would arise from replacing the known template response by its large-$L$
leading term. The resulting theorem concerns a prescribed fractional branch
with an exponentially small logarithmic-energy remainder. It neither derives
that remainder from quantum field theory nor covers arbitrary resummation,
unknown running, or general branch-cut spectral data.
